\documentclass[11pt]{article}
\usepackage{amsmath,amssymb,amsthm}
\usepackage[margin=1in]{geometry}
\newtheorem{theorem}{Theorem}
\newtheorem{lemma}{Lemma}
\newtheorem{proposition}{Proposition}
\newtheorem{corollary}{Corollary}
\theoremstyle{definition}
\newtheorem{remark}{Remark}
\newtheorem{observation}{Observation}
\newtheorem{conjecture}{Conjecture}

\newcommand{\Z}{\mathbb{Z}}
\newcommand{\C}{\mathbb{C}}
\newcommand{\R}{\mathbb{R}}
\newcommand{\N}{\mathbb{N}}
\newcommand{\Lam}{\Lambda}
\newcommand{\ip}[2]{\langle #1,\, #2\rangle}
\DeclareMathOperator{\supp}{supp}
\DeclareMathOperator{\Ind}{Ind}
\DeclareMathOperator{\ch}{ch}
\DeclareMathOperator{\mult}{mult}

\title{Two pruned proof routes for the $d{=}4$ fiber conjecture\\
$G_\lambda(i)=0\iff\lambda=(2,2)$}
\author{Clio}
\date{2026-06-05}

\begin{document}
\maketitle

\begin{abstract}
This is a negative-results note. Two imported-framework routes for the
full-support / $d{=}4$ fiber conjecture are investigated and found
\emph{not} to close it: (1) the Foulkes / $\omega$-twisted-plethysm
column-antisymmetrization mechanism, which controls only the
dominance-extreme (boundary) constituents and has no slot to act on an
\emph{ordinary} power $\psi^m$; and (2) the eigenvalue-multiplicity
framing, which fails on a type mismatch ($G_\lambda(i)$ is a Gaussian
integer, multiplicities are nonnegative reals) and on a structural
mismatch (one conjugacy class vs.\ a weighted sum over all square-types).
Each route is recorded precisely so it is not re-attempted. The portable
nugget extracted en route is the $\omega$-derivation of the conjugation
symmetry $G_{\lambda'}(i)=i^m\overline{G_\lambda(i)}$. A final synthesis
records that the live route is the author's own real-norm machinery.
\end{abstract}

\section*{Setup and the conjecture}

Let $\Lam$ be the ring of symmetric functions, with $h_r$ the complete
homogeneous and $e_r$ the elementary symmetric functions, and let
$\omega$ be the standard involution, $\omega(h_r)=e_r$. Fix
$i=\sqrt{-1}$ and set
\[
  \psi \;=\; h_2 + i\,e_2 \;\in\; \C\otimes\Lam .
\]
For a partition $\lambda\vdash 2m$ define the \emph{$d{=}4$ fiber value}
\[
  G_\lambda(i)\;=\;\ip{s_\lambda}{\psi^m},
\]
the Hall-inner-product coefficient of the Schur function $s_\lambda$ in
the degree-$2m$ symmetric function $\psi^m$. Equivalently, via the
character expansion of $h_2$ and $e_2$,
\[
  (1+i)^m\,G_\lambda(i)\;=\;\sum_{k=0}^m \binom{m}{k}\, i^k\,
  \chi^\lambda\!\bigl(2^{\,m-k}1^{\,2k}\bigr),
\]
a weighted sum of irreducible $S_{2m}$-characters over the involution
(``square-type'') classes $2^{m-k}1^{2k}$.

\begin{conjecture}[$d{=}4$ fiber]\label{conj:main}
For every $\lambda\vdash 2m$,
\[
  G_\lambda(i)=0 \iff \lambda=(2,2).
\]
\end{conjecture}

This is verified for all $\lambda\vdash 2m$ with $m\le 8$ (i.e.\ $n\le16$).
The open content is the \emph{full-support} claim: $\supp_\C(\psi^m)$ is
\emph{every} $\lambda\vdash 2m$ for $m\ge 3$, while for $m=2$ the only
missing shape is the self-conjugate rectangle $(2,2)$. (Here
$\supp_\C(\psi^m)=\{\lambda : G_\lambda(i)\ne0\}$.)

Two routes below were pursued to establish this support statement from an
external framework. Both are pruned.

\section{The Foulkes / $\omega$-twisted-plethysm route does not reach the interior}

We first record the one portable consequence of the $\omega$-involution,
then explain why the associated plethysm machinery does not transfer.

\begin{observation}[Conjugation symmetry, via $\omega$]\label{obs:conj}
Since $\omega(h_2)=e_2$ and $\omega(e_2)=h_2$,
\[
  \omega(\psi)=\omega(h_2+i\,e_2)=e_2+i\,h_2
  = i\,(h_2 - i\,e_2)=i\,\overline{\psi},
\]
where $\overline{\,\cdot\,}$ is complex conjugation of coefficients
(applied to a real basis of $\Lam$). As $\omega$ is a ring homomorphism
and conjugation is multiplicative on $\C\otimes\Lam$,
\[
  \omega(\psi^m)=\bigl(\omega(\psi)\bigr)^m
  =\bigl(i\,\overline{\psi}\bigr)^m = i^m\,\overline{\psi^m}.
\]
Pairing with $s_\lambda$ and using $\ip{s_\lambda}{\omega(f)}
=\ip{\omega(s_\lambda)}{f}=\ip{s_{\lambda'}}{f}$ (self-adjointness of
$\omega$, $\omega(s_\lambda)=s_{\lambda'}$) together with
$\ip{s_\lambda}{\overline{f}}=\overline{\ip{s_\lambda}{f}}$ for $s_\lambda$
real, we obtain
\[
  G_{\lambda'}(i)=\ip{s_{\lambda'}}{\psi^m}
  =\ip{s_\lambda}{\omega(\psi^m)}
  = i^m\,\overline{\ip{s_\lambda}{\psi^m}}
  = i^m\,\overline{G_\lambda(i)} .
\]
\end{observation}

This conjugation symmetry $G_{\lambda'}(i)=i^m\,\overline{G_\lambda(i)}$
was already known; the value of Observation~\ref{obs:conj} is that the
$\omega$-twist gives it in two lines, and it is the lens through which the
self-conjugate shape $(2,2)$ becomes the natural candidate vanisher
($\lambda=\lambda'$ forces $G_\lambda(i)=i^m\overline{G_\lambda(i)}$, a
phase constraint).

\paragraph{The lead.} The route was suggested by MathOverflow~\#501127
(answers of Wildon and Za\u{\i}mi) on the rectangle in a Foulkes-type
plethysm: $s_{(k^n)}$ occurs in $h_n[h_k]$ iff $k$ is even. Two proofs
are given there.

\begin{remark}[Wildon's column-antisymmetrization]\label{rem:wildon}
Wildon's argument tracks a \emph{single} column-antisymmetrized
polytabloid through the surjection $M^{(k^n)}\twoheadrightarrow H^{(k^n)}$
onto the Foulkes module. Antisymmetrizing one column produces a
$(-1)^k$ block-swap sign; for $k$ odd the image vanishes, for $k$ even it
survives. This decides occurrence of the rectangle $s_{(k^n)}$, which is
the \emph{dominance-minimal} constituent of $h_n[h_k]$.
\end{remark}

\begin{remark}[Za\u{\i}mi's $\omega$-twist]\label{rem:zaimi}
Za\u{\i}mi uses $\omega(s_\lambda[s_\mu])=s_\lambda[s_{\mu'}]$ for
$|\mu|$ even, reducing the minimal-constituent question, after the twist,
to a lex-leading-monomial computation. This is again an argument about a
single extremal constituent (now lex-leading after the twist), not the
interior.
\end{remark}

\begin{proposition}[Why the route is pruned]\label{prop:foulkes-pruned}
The Foulkes column-antisymmetrization mechanism can decide the
dominance-extreme constituents of $\supp(\psi^m)$ (the boundary of the
support polytope) but is not a full-support engine for
Conjecture~\ref{conj:main}. Two obstructions are decisive.
\end{proposition}

\begin{proof}[Discussion]
\emph{(i) Extremal, not interior.} Both
Remarks~\ref{rem:wildon}--\ref{rem:zaimi} track \emph{one} extremal
constituent through \emph{one} surjection. The $(-1)^k$ cancellation is a
statement about the dominance-minimal (resp.\ lex-leading) polytabloid;
it gives no control of an interior $\lambda$, for which one would need a
per-$\lambda$ cancellation analysis that is genuinely shape-specific.
Indicative of this, the minimal-constituent classification of
Paget--Wildon (arXiv:1007.2946) is complete only for $n\le 5$; there is no
uniform interior statement to import.

\emph{(ii) No plethysm slot for the sign to act in.} Both arguments
require a genuine \emph{plethysm} / wreath-product induction,
\[
  h_m[h_2]=\ch\,\Ind_{S_2\wr S_m}^{S_{2m}} \mathbf{1},
\]
so that the $(-1)^k$ block-swap acts in the passage from the
permutation module $M^{(2^m)}$ down to the Foulkes quotient $H^{(2^m)}$.
But $\psi^m=(h_2+i\,e_2)^m$ is an \emph{ordinary product} --- a power of a
fixed degree-$2$ element of $\C\otimes\Lam$ --- and lives in the
$M^{(2^m)}$ world \emph{before} any quotient to a Foulkes module. There is
no plethystic bracket, hence no surjection onto $H^{(2^m)}$ in which the
sign mechanism has a slot to act. (The $i\,e_2$ summand only deforms the
ordinary product; it does not manufacture a wreath induction.)
Consequently the precise mechanism that proves the rectangle theorem has
no analogue here at the interior.
\end{proof}

\noindent\textbf{Conclusion of \S1.} Foulkes/$\omega$-plethysm decides
the boundary of $\supp(\psi^m)$ only. It does not establish full support,
and is pruned as a route to Conjecture~\ref{conj:main}. The single
retained item is the $\omega$-derivation of
Observation~\ref{obs:conj}.

\section{$G_\lambda(i)$ is not an eigenvalue multiplicity}

\paragraph{The lead (GR-1).} For an order-$4$ element $\sigma\in S_{2m}$
and the irreducible representation $\rho_\lambda$, the multiplicity of the
eigenvalue $i$ in $\rho_\lambda(\sigma)$ is
\[
  \mult_i(\lambda,\sigma)=\frac14\sum_{j=0}^{3} i^{-j}\,
  \chi^\lambda(\sigma^j).
\]
The lead asked whether $G_\lambda(i)$ equals, up to a fixed scalar,
$\mult_i(\lambda,\sigma)$ for some natural $\sigma$. The relevant
references are Staroletov (arXiv:2501.17571), on the spectrum of
$\rho_\lambda(\sigma)$, and Amrutha--Prasad--Velmurugan
(arXiv:2308.08146), relating an invariant vector to $1$ being an
eigenvalue.

\begin{proposition}[Why the route is pruned]\label{prop:eig-pruned}
$G_\lambda(i)$ is not (up to a fixed scalar) an eigenvalue multiplicity
$\mult_i(\lambda,\sigma)$ of any order-$4$ element, and neither cited
theorem short-cuts Conjecture~\ref{conj:main}.
\end{proposition}

\begin{proof}[Discussion]
\emph{(i) Type mismatch.} For every $\lambda$ and every $\sigma$,
$\mult_i(\lambda,\sigma)\in\Z_{\ge0}$ is a nonnegative integer. By
contrast $G_\lambda(i)$ is a genuine Gaussian integer that is
\emph{generically complex} (nonreal). A nonnegative real cannot equal a
generically nonreal quantity up to a fixed scalar, so no identification
$G_\lambda(i)=c\cdot\mult_i(\lambda,\sigma)$ can hold with $c$ independent
of $\lambda$.

\emph{(ii) Structural mismatch.} The multiplicity $\mult_i$ is supported
on the four powers $\{\sigma^0,\sigma^1,\sigma^2,\sigma^3\}$ of one
conjugacy class and uses the $4$-cycle character values
$\chi^\lambda(4,\dots)$. The fiber value $G_\lambda(i)$, on the other
hand, is a $\binom{m}{k}i^k$-weighted sum over \emph{all} square-types
$2^{m-k}1^{2k}$ with an imaginary sign-twist: it is the projection of
$s_\lambda$ against the \emph{fixed symmetric function} $\psi^m$, not the
class indicator of a single conjugacy class. These are different linear
functionals on the character table.

\emph{(iii) Computation.} For all $\lambda\vdash 2m$ with $m\le 8$ and
every natural order-$4$ element $\sigma$, $G_\lambda(i)$ is not
proportional to $\mult_i(\lambda,\sigma)$; the zero-sets do not even
coincide. The multiplicity $\mult_i$ vanishes only at the two
one-dimensional representations, whereas $G_\lambda(i)$ vanishes only at
$(2,2)$.

\emph{(iv) The cited theorems do not help.} Staroletov's result gives the
\emph{spectrum} of $\rho_\lambda(\sigma)$ (which roots of unity occur),
not the multiplicities; it confirms that $i$ occurs generically but says
nothing about the weighted count $G_\lambda(i)$. The Amrutha--Prasad--%
Velmurugan exceptional pair (relating $\{(2,2),(3,1)\}$ to an invariant
vector) concerns an order-$3$ element and the eigenvalue $1$, not an
order-$4$ element and $i$.
\end{proof}

\noindent\textbf{Conclusion of \S2.} The eigenvalue-multiplicity framing
does not capture the invariant $G_\lambda(i)$, and is pruned. Neither
Staroletov nor Amrutha--Prasad--Velmurugan short-cuts the $(2,2)$
conjecture.

\section{Synthesis: the proof must come from the norm/M\"obius machinery}

The imported-framework routes now pruned for the $d{=}4$ fiber include:
CSP carriers; Albion-type $z$-asymmetric (root-of-unity-twisted)
plethysm; the Hermitian family $\Xi=h_2^2+e_2^2$; the Bethe/spectral
home; Foulkes support (\S1); and eigenvalue multiplicity (\S2). None is a
full-support engine.

The \emph{live} route is the author's own hand-built machinery. For two
rows,
\[
  G_{(a,b)}(i)=r_b-r_{b-1},\qquad
  r_\ell=[u^\ell]\,\bigl(u^2+(1+i)u+1\bigr)^m ,
\]
and the real norm $N_\ell=|r_\ell|^2$ is strictly increasing in $\ell$
with the \emph{single} tie $N_1=N_2$ occurring only at $m=2$ (the shape
$(2,2)$). This reduces the two-row instance of
Conjecture~\ref{conj:main} to one finite inequality --- ``Lemma~1'', a
M\"obius-orbit / norm-ratio bound. The standing structural question is
whether this real-norm proxy lifts from two rows to general $\lambda$: one
seeks a Gram-determinant analogue of $N_\ell$ in which the unique
coincidence again pins down the self-conjugate rectangle $(2,2)$. That is
the direction recorded as live.

\end{document}
